\section{Methods}
\subsection{Activation Screening by Linear Programming}
\label{sec:feasible-space}

The low-activation feasible space is defined over the eleven-element
\ch{Cr-Fe-Hf-Mo-Nb-Re-Ta-Ti-V-W-Zr} system, sampled on a \(5\)~at.\% grid, using
the \emph{FLARE} activation model~\cite{Marquez2026FLARE} together with linear programming.
\emph{FLARE} is a reduced-order model that estimates the neutron-activation
response of an alloy, gas generation, displacement damage, decay heat, and contact
dose rate, as a composition-weighted sum of precomputed elemental contributions,
providing a quantitative basis for assessing irradiation-induced phenomena across
the design space.

The responses are evaluated under the irradiation conditions summarized in
\Cref{tab:ind_variables}, a two-year irradiation time with decay times of one
second, one month, and 100 years, corresponding to the immediate post-irradiation,
maintenance, and long-term waste management scenarios. Under these conditions, the
responses of pure W and EUROFER97 serve as reference baselines for plasma-facing
and structural materials, respectively, and a composition is retained if each of
its responses stays within a deviation factor of \(\phi=10\) of the corresponding
baseline value.

\begin{table}[H]
\centering
\caption{Summary of irradiation condition inputs used in the \emph{FLARE} model to
evaluate activation responses for candidate alloys under first-wall neutron
environments.}
\label{tab:ind_variables}
\begin{tabularx}{\columnwidth}{lX}
\toprule
\multicolumn{2}{c}{\textbf{Independent Variables}} \\
\midrule
First Wall Average Neutron Flux & \(8\times10^{13}~\mathrm{n/cm^2 s}\) \\
Irradiation Time               & 2 years \\
Decay Time                     & 1 second \& 1 month \& 100 years \\
\bottomrule
\end{tabularx}
\end{table}

The conventional way to apply these baselines is to generate every composition,
evaluate its responses, and discard those that fall outside the limits. For a
\(5\)~at.\% grid this requires testing on the order of \(10^{7}\) compositions one
by one, which is slow and scales poorly as elements or resolution are added. We
avoid this entirely by using the structure of the screening problem itself, since
the FLARE responses are linear in the alloy composition
\(\mathbf{x}=(x_1,\dots,x_N)\), directly for the atomic-fraction-averaged
displacement damage and, for the mass-weighted responses, once each criterion
is multiplied through by their common alloy-mass denominator. Linearity is what makes the shortcut
possible, since when a response varies linearly with composition, the boundary at which it crosses
a baseline limit is a flat surface in composition space, so each limit simply
divides the space into an allowed and a disallowed side rather than having to be
tested point by point. Each baseline limit is therefore a linear inequality in
\(\mathbf{x}\), and stacking all of them, with the simplex condition
\(\sum_i x_i = 1\) and the bounds \(0 \le x_i \le 1\), gives a single linear
system \(\mathbf{G}\mathbf{x} \le \mathbf{h}\). The compositions satisfying
every limit are exactly those lying inside the region enclosed by these surfaces, a
convex polytope

\begin{equation}
P=\{\mathbf{x}\in\Delta : \mathbf{G}\mathbf{x}\le\mathbf{h}\}.
\label{eq:polytope}
\end{equation}

The feasible region is thus obtained directly as \(P\) by linear programming, and
the \(5\)~at.\% grid compositions inside it are enumerated without evaluating the
infeasible majority of the grid. This is substantially faster than
generate-and-filter and, importantly, it is not specific to FLARE, as the same
construction applies to any screening model whose responses are linear in
composition, FLARE being one such model. The resulting low-activation feasible
space is the starting point for the property evaluation and graded-path design that
follow.

\subsection{Property Calculations}
\label{sec:feasible-space-props}

Following activation screening, the surviving compositions were characterized by
high-throughput calculations of the properties on which the design depends. Because
direct experimental characterization is impractical for the number of candidates
considered, physically motivated models were used as first-order estimates: the
Maresca--Curtin model for yield strength, the Pugh ratio for intrinsic ductility,
and CALPHAD for phase stability, melting behavior, and thermophysical properties.

Yield strength was estimated with the Maresca--Curtin (MC) model for
body-centered cubic random solid solutions~\cite{maresca2020mechanistic},
following the methodology previously applied to the screening of refractory
multi-principal-element alloys~\cite{vela2023data}. The model relates strength to
the energetics of thermally activated dislocation glide, with the
temperature-dependent yield strength given by
\begin{equation}
\tau_y(T,\dot{\varepsilon})
= \tau_{y0}
\exp\!\left[
-\frac{1}{0.55}
\left(
\frac{kT}{\Delta E_b}
\ln\!\left(\frac{\dot{\varepsilon}_0}{\dot{\varepsilon}}\right)
\right)^{0.91}
\right],
\label{eq:yield_stress}
\end{equation}
where $\tau_{y0}$ is the athermal yield strength, $\Delta E_b$ is the
activation barrier for dislocation glide, $T$ is the absolute temperature,
$\dot{\varepsilon}$ is the applied strain rate, $\dot{\varepsilon}_0$ is the
reference strain rate, and $k$ is Boltzmann's constant, with the tensile
yield strength obtained through the appropriate Taylor-factor conversion. The
physical inputs, including atomic volumes, elastic constants ($C_{11}$, $C_{12}$,
and $C_{44}$), shear moduli, and Poisson ratios, were assembled from literature
data and combined by the rule of mixtures. Yield strengths were evaluated at 600
and $650^{\circ}\mathrm{C}$, the service temperatures considered for the
plasma-facing and structural sides of the wall, at a strain rate of
$10^{-3}\,\mathrm{s}^{-1}$, and are used as design criteria in
Section~\ref{sec:joinable}.

Intrinsic ductility was assessed with the Pugh ratio \(B/G\), the ratio of bulk to
shear modulus, in the modern convention in which larger values indicate greater
ductility~\cite{pugh1954,singh2023ductility}. The elemental Poisson ratios
\(\nu_i\) were combined by the rule of mixtures, \(\bar{\nu}=\sum_i x_i \nu_i\),
where $x_i$ is the atomic fraction of element $i$, and through the isotropic
relations \(E=3B(1-2\nu)\) and \(G=E/2(1+\nu)\) the descriptor reduces to
\begin{equation}
\frac{B}{G}=\frac{2\,(1+\bar{\nu})}{3\,(1-2\bar{\nu})}.
\label{eq:pugh}
\end{equation}
A composition is considered intrinsically ductile when \(B/G \ge 2.6\), a stricter
cutoff than the conventional \(B/G > 1.75\) criterion~\cite{singh2023ductility}.
As a rule-of-mixtures prior, this descriptor serves as a screening-level estimate.

Phase stability, melting behavior, and thermophysical properties were obtained with
the CALPHAD (CALculation of PHAse Diagrams) method, in which the Gibbs free energy
of every candidate phase is described by parametric models assessed from
experimental and first-principles data and extrapolated to multicomponent
alloys~\cite{li2023calphad,conway2022thermodynamic}. Calculations were performed
with Thermo-Calc's TC-Python API and the TCHEA6 database, assessed for
high-entropy and refractory multi-principal-element systems. Equilibrium phase
fractions, including the BCC solid solution and competing intermetallic phases such
as sigma, Laves, mu, and chi, were computed at 500, 600, and
$650^{\circ}\mathrm{C}$. The solidus temperature \(T_s\) was obtained with the
Liquidus and Solidus Temperature property model, and the thermal conductivity at
$25^{\circ}\mathrm{C}$ and the linear coefficient of thermal expansion (CTE)
\(\alpha\) at 600 and $650^{\circ}\mathrm{C}$ with the Equilibrium with Freeze-in
temperature property model.

\subsection{Graded-Path Design}
\label{sec:joinable}

The activation-feasible compositions were first screened for single-phase BCC
stability and heat removal, retaining a composition only if
\begin{equation}
f_{\mathrm{BCC}}(T) \ge 0.98 \quad (T = 500,\ 600,\ 650\ ^{\circ}\mathrm{C}),
\qquad \kappa_{25\,^{\circ}\mathrm{C}} > 16~\mathrm{W\,m^{-1}\,K^{-1}},
\label{eq:filters}
\end{equation}
where \(f_{\mathrm{BCC}}\) is the BCC phase fraction and \(\kappa\) the thermal
conductivity. The surviving compositions were organized into a simplex-lattice
graph~\cite{nimplex}, in which two compositions \(u\) and \(v\) are neighbors when
they differ by a single grid step,
\(\mathbf{x}_v-\mathbf{x}_u = 0.05\,(\mathbf{e}_b-\mathbf{e}_a)\), one element
decreased and another increased by one increment. Graph representations of the
design space have likewise been used for the design of compositionally graded
structures~\cite{multiterminalCGA2024}. Because the screens remove compositions
unevenly, the graph separates into connected subgraphs, and the largest is carried
forward as the design region since it offers the widest contiguous space for a
graded transition~\cite{obrist2026graph}.

Each composition of the design region is classified by the property the two sides
of the first wall require. A composition is a plasma-facing candidate,
Material~A, or a structural candidate, Material~B, when
\begin{equation}
\text{A:}\ \ T_s > 2200~^{\circ}\mathrm{C}\ \text{and}\ \sigma_y^{600} \ge 300~\mathrm{MPa},
\qquad
\text{B:}\ \ B/G \ge 2.6\ \text{and}\ \sigma_y^{650} \ge 250~\mathrm{MPa},
\label{eq:classes}
\end{equation}
where \(\sigma_y^{T}\) is the yield strength at \(T\)~\(^{\circ}\)C. The strength
limits are set at the level of the reference structural steel
EUROFER97~\cite{lindau2005eurofer}, so that every candidate is at least as strong as
the steel it would replace at the service temperatures.
Compositions satisfying both definitions form Material~AB, and those satisfying
neither are classified as Intermediate. The two ends of the gradient are chosen by
the function of each side of the wall, as in previous gradient design
studies~\cite{kirk2021monotonic}: the plasma-facing (PFM) end is the
highest-solidus composition of Material~A or AB, and the structural end the
highest-Pugh-ratio composition of Material~B or AB,
\begin{equation}
v_{\mathrm{PFM}} = \arg\max_{v \in A \cup AB} T_s(v), \qquad
v_{\mathrm{str}} = \arg\max_{v \in B \cup AB} (B/G)(v).
\label{eq:ends}
\end{equation}

A graded path \(P=(v_0,v_1,\dots,v_L)\), with \(v_0=v_{\mathrm{PFM}}\) and
\(v_L=v_{\mathrm{str}}\), is a sequence of neighboring compositions, each step
corresponding to one deposited layer. Every layer must be a design material that is
at least as strong as EUROFER97 at both service temperatures and a multicomponent
alloy, so the admissible compositions are
\begin{equation}
\mathcal{V} = \{\, v \in A \cup B \cup AB :\ \sigma_y^{600}(v) \ge 300~\mathrm{MPa},\
\sigma_y^{650}(v) \ge 250~\mathrm{MPa},\ N_{\mathrm{el}}(v) \ge 3 \,\},
\label{eq:admissible}
\end{equation}
where \(N_{\mathrm{el}}\) is the number of constituent elements. Two monotonicity
constraints are imposed on every step \(u \to v\),
\begin{equation}
\alpha_T(v) \ge \alpha_T(u) \quad (T = 600,\ 650\ ^{\circ}\mathrm{C}),
\qquad T_s(v) \le T_s(u).
\label{eq:monotone}
\end{equation}
A monotonic CTE keeps the thermal-expansion mismatch at every interface in the same
direction, from the low-expansion plasma-facing alloy to the higher-expansion
structural alloy, so that no interlayer reverses the strain gradient across the
wall~\cite{kirk2021monotonic}. A monotonic solidus follows the temperature profile
of the wall, which falls from the plasma-facing surface toward the coolant, so that
no layer has a lower melting point than a layer behind it, in the same way that a
monotonic liquidus prevents remelting of previously deposited layers during graded
deposition~\cite{glerum2026ss316tzm}. These constraints are imposed on the graph
itself: each undirected edge is expanded into its two directed steps, and a step is
kept only if both compositions belong to \(\mathcal{V}\) and it satisfies
\Cref{eq:monotone}.

The path cost follows Kirk et al.~\cite{kirk2021monotonic}, a monotonicity
penalty plus a small path-length weight,
\begin{equation}
c(P) = \mathrm{LOI}(P) + w\,l(P), \qquad w \ll 1,
\label{eq:cost}
\end{equation}
where the loss of invariance, LOI, is the total change of the constrained
properties against their required direction and \(l\) is the path length in
composition space,
\begin{equation}
\mathrm{LOI}(P)=\sum_{k=1}^{L}\Big[\sum_{T}\max\!\big(\alpha_T(v_{k-1})-\alpha_T(v_k),0\big)
+\max\!\big(T_s(v_k)-T_s(v_{k-1}),0\big)\Big],
\qquad
l(P)=\sum_{k=1}^{L}\lVert \mathbf{x}_{v_k}-\mathbf{x}_{v_{k-1}}\rVert_2 .
\label{eq:loi}
\end{equation}
Because the monotonicity constraints are imposed as hard rules, \(\mathrm{LOI}=0\) on
every admissible path, and since every step is a single 5~at.\% exchange,
\(l = 0.05\sqrt{2}\,L\). The cost therefore reduces to the number of steps, and the
design path is the shortest monotonic path,
\begin{equation}
P^{*} = \arg\min_{P}\, L(P),
\label{eq:shortest}
\end{equation}
which is found exactly by Dijkstra's algorithm~\cite{dijkstra1959note} with unit
step weights on the directed graph of admissible monotonic steps. Because every
composition property is computed once and stored with the graph, the search
returns the optimal path exactly rather than approximately, in contrast to the
sampling-based planners used previously~\cite{kirk2021monotonic,kirk2018gradient}.
Several paths can share the minimum length; they are equally optimal under this
cost, and the algorithm returns one of them deterministically. The returned path is
verified step by step against \Cref{eq:admissible,eq:monotone}.
